\documentclass[11pt]{article}
\usepackage[margin=1in]{geometry}
\usepackage[T1]{fontenc}
\usepackage{lmodern,microtype,hyperref,enumitem,amsmath,amssymb,xcolor,booktabs}
\usepackage[numbers,sort&compress]{natbib}
\hypersetup{colorlinks=true,urlcolor=blue,citecolor=blue,linkcolor=blue}
\title{Q25: Comprehension Check}
\author{DLS 2026 --- Kukla}
\date{October 2026}
\begin{document}
\maketitle
\textbf{Comprehension check.} A team uses model A to generate training data \emph{and} model A (same model)
to judge which generated answers are best, then fine-tunes model B on the winners and reports that B
improved---as scored by model A. Name the specific risk in this setup and one change that would make the
reported improvement more trustworthy.

\vspace{0.8cm}\hrule\vspace{4pt}
{\small\textbf{Answer.} The risk is \textbf{circularity / self-preference bias} (a form of reward hacking):
because the same model generates the data and judges quality, its own stylistic biases are baked into both
the training signal and the evaluation, so B can score well on A's judgments without genuinely improving on
what humans care about---an illusory gain (and, over iterations, a path toward model collapse). A
trustworthy change: evaluate B on an \emph{independent}, human-labeled held-out set (or at least a different
judge model and fresh real data), decoupling the generator from the evaluator.}
\end{document}
